% Shared typesetting template. `/autonovel:typeset` copies this file
% into `books/<name>/typeset/novel.tex` and substitutes the three
% placeholders @TITLE@, @AUTHOR@, @SERIES_NAME@ at build time. Hand-
% edit the copy, never this template, for per-book tweaks.

\documentclass[11pt, openany]{book}

% === GEOMETRY: Trade paperback (5.5 x 8.5 inches) ===
\usepackage[
  paperwidth=5.5in,
  paperheight=8.5in,
  inner=0.85in,
  outer=0.65in,
  top=0.75in,
  bottom=0.85in,
  headheight=14pt
]{geometry}

% === FONTS ===
\usepackage{fontspec}
% Primary font: EB Garamond. If fontconfig can't find it (common on
% Chromebook / WSL / fresh macOS), fontspec walks its fallback chain
% trying many similar names — the user sees a "stepping through many
% fonts by name" failure with no clear root cause. We catch the
% missing-font case explicitly via \IfFontExistsTF and fall back to a
% widely-available serif so typeset still produces a readable PDF.
% Tier of preference:
%   1. EB Garamond     (the design choice; Google Fonts free)
%   2. Garamond        (system Garamond if present — macOS ships one)
%   3. TeX Gyre Pagella (always present in tectonic's bundled set;
%                        Palatino-shaped, similar feel to EB Garamond)
%
% To install EB Garamond:
%   Linux/WSL: sudo apt-get install -y fonts-ebgaramond && fc-cache -fv
%   macOS:     brew install --cask font-eb-garamond
%   Or:        autonovel install-export-tools --apply  (handles both)
\IfFontExistsTF{EB Garamond}{
  \setmainfont{EB Garamond}[
    Ligatures=TeX,
    Numbers=OldStyle
  ]
}{%
  \IfFontExistsTF{Garamond}{
    \setmainfont{Garamond}[Ligatures=TeX]
  }{%
    % TeX Gyre Pagella is bundled with tectonic and always available.
    % Slightly different metrics than EB Garamond (wider em; less
    % calligraphic italic) but type-page-perfect and free.
    \setmainfont{TeX Gyre Pagella}[Ligatures=TeX]
    \typeout{*** WARNING: EB Garamond not installed; using TeX Gyre Pagella fallback.}
    \typeout{*** Install EB Garamond for the design-intended look:}
    \typeout{***   sudo apt-get install -y fonts-ebgaramond  (Linux/WSL)}
    \typeout{***   brew install --cask font-eb-garamond       (macOS)}
    \typeout{***   autonovel install-export-tools --apply}
  }
}

% === TYPOGRAPHY ===
\usepackage{microtype}
\usepackage{setspace}
\setstretch{1.12}
\usepackage{parskip}
\setlength{\parindent}{1.5em}
\setlength{\parskip}{0pt}

% === WIDOW / ORPHAN / HYPHEN PENALTIES ===
% User 2026-04-30 reported "words get split across pages, perhaps
% headings on the last lines of a page". LaTeX's defaults are
% permissive about widows (a single line of a paragraph at the top
% of a page) and orphans (last line of a paragraph alone at the
% bottom). Cranking the penalties forces TeX to push widow/orphan
% lines to the next page — costs a few half-empty pages but
% drastically improves the perceived quality. 10000 is "infinitely
% bad" in TeX's penalty units; the convention for novel typesetting
% is to use 10000 for both.
\widowpenalty=10000
\clubpenalty=10000
% Don't break a hyphenated word across pages (very ugly).
\brokenpenalty=10000
% Discourage starting a page with a hyphenation continuation.
\interlinepenalty=100
% Allow a slight bottom-rag (lines NOT forced to flush against the
% bottom of every page) so widow/orphan adjustments don't produce
% large rubber-spaced text. \raggedbottom is the convention for
% novel pages where consistent leading is more important than
% aligned bottom margins.
\raggedbottom

% === HEADING-ON-LAST-LINE PROTECTION ===
% titlesec's \chapter / \section etc. don't always respect needspace
% — a section heading can land at the bottom of a page with its
% first prose paragraph on the next page. \needspace below each
% heading forces a minimum of 4 lines of room; if not available,
% TeX advances to the next page first.
\usepackage{needspace}

% === GRAPHICS ===
\usepackage{graphicx}

% \cleartoverso advances to the next even (verso) page, adding a
% blank page if needed. Used by mechanical/latex.py's plate
% renderer to put before-chapter plates on the verso so the
% chapter heading lands on the facing recto — the published-book
% convention for full-page chapter-opening images.
\newcommand{\cleartoverso}{%
  \clearpage%
  \ifodd\value{page}%
    \hbox{}%
    \thispagestyle{empty}%
    \clearpage%
  \fi%
}

% === DROP CAPS ===
\usepackage{lettrine}
\setcounter{DefaultLines}{2}

% === HEADERS AND FOOTERS ===
% Classic novel typography: book title on the verso (left-even), the
% chapter mark on the recto (right-odd). The recto deliberately uses
% `Chapter <Roman>` rather than `\leftmark` (chapter title) — chapter
% titles vary in length, can carry italics or special characters, and
% in author-testing 2026-04-25 the previous `\textit{\leftmark}` setup
% bled italic chapter-title fragments into the running header in
% surprising ways. The Roman chapter number is short, stable, and
% always renders cleanly.
\usepackage{fancyhdr}
\pagestyle{fancy}
\fancyhf{}
% Running header: book title on verso, chapter mark on recto. The
% recto reads `\rightmark` (set by `\chaptermark` for numbered
% chapters and by manual `\markboth` for unnumbered chapters like
% Preface / Introduction / Glossary / Appendix). Without the
% `\rightmark` indirection, the appendix's running header would
% show "Chapter 24" — the last \thechapter value — because
% `\chapter*` doesn't auto-update the header. User 2026-04-30 hit
% exactly that bug.
\fancyhead[LE]{\small\textsc{@TITLE@}}
\fancyhead[RO]{\small\textsc{\rightmark}}
\fancyfoot[C]{\thepage}
\renewcommand{\headrulewidth}{0pt}

\fancypagestyle{plain}{
  \fancyhf{}
  \fancyfoot[C]{\thepage}
  \renewcommand{\headrulewidth}{0pt}
}

% === CHAPTER STYLE ===
\usepackage{titlesec}

\newcommand{\scenebreak}{%
  \par\vspace{0.6\baselineskip}%
  \noindent\hfil%
  \IfFileExists{../art/pdf/scene_break.pdf}{%
    \includegraphics[width=1.2in]{../art/pdf/scene_break.pdf}%
  }{%
  \IfFileExists{../art/scene_break.png}{%
    \includegraphics[width=1.2in]{../art/scene_break.png}%
  }{%
    {\small\symbol{"2022}\quad\symbol{"2022}\quad\symbol{"2022}}%
  }}%
  \hfil%
  \par\vspace{0.6\baselineskip}%
}

\renewcommand{\thechapter}{\Roman{chapter}}
\titleformat{\chapter}[display]
  {\normalfont\centering}
  {\vspace*{1.5in}\footnotesize\textsc{chapter \thechapter}}
  {4pt}
  {\Large\itshape}
  [\vspace{8pt}{\small--- $\diamond$ ---}\vspace{0.5in}]

\titlespacing*{\chapter}{0pt}{0pt}{0pt}

% Lowercase chapter name in header
% \chaptermark fires for numbered \chapter; we set the right mark
% to "Chapter <Roman>" so the recto running header reads
% "Chapter VII" via \rightmark in the fancyhead config above.
% Unnumbered chapter*'s manually call \markboth with their name
% (see back_matter.py / front_matter.py).
\renewcommand{\chaptermark}[1]{\markboth{}{Chapter \thechapter}}

% === TITLE PAGE DESIGN ===
\newcommand{\makenoveltitle}{%
  \thispagestyle{empty}
  \begin{center}
  \vspace*{2in}

  {\Huge\textsc{@TITLE@}}\\[0.4in]

  {\small------\quad$\diamond$\quad------}\\[0.5in]

  {\large\textit{A Novel}}\\[1in]

  {\Large\textsc{@AUTHOR@}}\\[1.5in]

  \end{center}
  \clearpage
}

% === HALF TITLE ===
\newcommand{\makehalftitle}{%
  \thispagestyle{empty}
  \vspace*{3in}
  \begin{center}
  {\large\textsc{@TITLE@}}
  \end{center}
  \clearpage
}

% === PDF METADATA ===
\usepackage{hyperref}
\hypersetup{
  pdftitle={@TITLE@},
  pdfauthor={@AUTHOR@},
  pdfsubject={@SERIES_NAME@},
  hidelinks
}

% === BEGIN DOCUMENT ===
\begin{document}

\frontmatter

% Full-bleed cover image (if exists). Prefers `cover_titled.png`
% (title + author overlaid by `/autonovel:cover-composite`) over the
% bare `cover.png` — the user 2026-04-30 hit "PDF missing the title
% overlay on the front cover image" because the template was reading
% the unoverlaid file. Both checks are guarded so a missing cover
% silently skips this section (the half-title still renders the
% title in clean type).
\IfFileExists{../art/cover_titled.png}{%
  \thispagestyle{empty}
  \newgeometry{margin=0pt}
  \noindent\includegraphics[width=\paperwidth, height=\paperheight, keepaspectratio]{../art/cover_titled.png}
  \restoregeometry
  \clearpage
}{%
  \IfFileExists{../art/cover.png}{%
    \thispagestyle{empty}
    \newgeometry{margin=0pt}
    \noindent\includegraphics[width=\paperwidth, height=\paperheight, keepaspectratio]{../art/cover.png}
    \restoregeometry
    \clearpage
    \typeout{*** WARNING: cover.png exists but cover_titled.png does not.}
    \typeout{*** Run /autonovel:cover-composite --book <name> to overlay}
    \typeout{*** title + author. The bare cover.png is used as a fallback}
    \typeout{*** but readers won't see the book title on the cover.}
  }{}%
}

% Half title
\makehalftitle

% Blank verso
\thispagestyle{empty}
\mbox{}
\clearpage

% Title page
\makenoveltitle

% Copyright / colophon
\thispagestyle{empty}
\vspace*{\fill}
\begin{center}

{\small This is a work of fiction.}\\[18pt]

{\small Typeset by autonovel.}

\end{center}
\vspace*{\fill}
\clearpage

% === FRONT MATTER (preface + introduction, optional) ===
% typeset writes `front_matter.tex` from `preface.md` +
% `introduction.md` if either exists in the book's root. Preface goes
% first (author voice), then introduction (typically more about the
% work itself). Both render as `\chapter*` so they're unnumbered and
% kept out of the chapter sequence, but added to the TOC.
\IfFileExists{front_matter.tex}{\input{front_matter.tex}}{}

% === TABLE OF CONTENTS ===
% User 2026-04-30 reported "the chapter names still don't show in
% toc". Real root cause: there was no \tableofcontents call AT ALL
% in the template — the document had no TOC. \chapter{<title>}
% auto-emits an \addcontentsline entry, but without
% \tableofcontents nothing renders it. Sits in \frontmatter zone
% so it's unnumbered and uses Roman page numbering.
\tableofcontents
\clearpage

\mainmatter

% === CHAPTERS ===
\input{chapters_content.tex}

% === END MATTER ===
\backmatter

% Optional appendix (real-world events, bios, sources, etc.). Written
% by `autonovel mechanical build-back-matter-tex` from `appendix.md`
% if the file exists in the book's root. The `\IfFileExists` guard
% silently skips inclusion when no appendix has been authored.
\IfFileExists{back_matter.tex}{\input{back_matter.tex}}{}

\thispagestyle{empty}
\vspace*{3in}
\begin{center}
{\small------\quad$\diamond$\quad------}
\end{center}

% === BACK COVER IMAGE (optional) ===
% Full-bleed back-cover image, parallel to the front-cover block in
% \frontmatter. User 2026-04-30: "there is a place for a back page
% image that we are not generating, need a way to do that." Drop a
% PNG at books/<book>/art/back_cover.png (compose via
% /autonovel:cover-back when shipped, or hand-author) and it
% renders here.
\IfFileExists{../art/back_cover.png}{%
  \clearpage
  \thispagestyle{empty}
  \newgeometry{margin=0pt}
  \noindent\includegraphics[width=\paperwidth, height=\paperheight, keepaspectratio]{../art/back_cover.png}
  \restoregeometry
  \clearpage
}{}

\end{document}
